\documentclass{amsart}
% For dealing with references we use the comment environment
\usepackage{verbatim}
\newenvironment{reference}{\comment}{\endcomment}
%\newenvironment{reference}{}{}
\newenvironment{slogan}{\comment}{\endcomment}
\newenvironment{history}{\comment}{\endcomment}

% For commutative diagrams we use Xy-pic
\usepackage[all]{xy}

% We use 2cell for 2-commutative diagrams.
\xyoption{2cell}
\UseAllTwocells

% We use multicol for the list of chapters between chapters
\usepackage{multicol}

% This is generally recommended for better output
\usepackage{lmodern}
\usepackage[T1]{fontenc}

% For cross-file-references

% Package for hypertext links:
\usepackage{hyperref}

% For any local file, say "hello.tex" you want to link to please

% Theorem environments.
%
\theoremstyle{plain}
\newtheorem{theorem}[subsection]{Theorem}
\newtheorem{proposition}[subsection]{Proposition}
\newtheorem{lemma}[subsection]{Lemma}

\theoremstyle{definition}
\newtheorem{definition}[subsection]{Definition}
\newtheorem{example}[subsection]{Example}
\newtheorem{exercise}[subsection]{Exercise}
\newtheorem{situation}[subsection]{Situation}

\theoremstyle{remark}
\newtheorem{remark}[subsection]{Remark}
\newtheorem{remarks}[subsection]{Remarks}

\numberwithin{equation}{subsection}

% Macros
%
\def\lim{\mathop{\mathrm{lim}}\nolimits}
\def\colim{\mathop{\mathrm{colim}}\nolimits}
\def\Spec{\mathop{\mathrm{Spec}}}
\def\Hom{\mathop{\mathrm{Hom}}\nolimits}
\def\Ext{\mathop{\mathrm{Ext}}\nolimits}
\def\SheafHom{\mathop{\mathcal{H}\!\mathit{om}}\nolimits}
\def\SheafExt{\mathop{\mathcal{E}\!\mathit{xt}}\nolimits}
\def\Sch{\mathit{Sch}}
\def\Mor{\mathop{\mathrm{Mor}}\nolimits}
\def\Ob{\mathop{\mathrm{Ob}}\nolimits}
\def\Sh{\mathop{\mathit{Sh}}\nolimits}
\def\NL{\mathop{N\!L}\nolimits}
\def\CH{\mathop{\mathrm{CH}}\nolimits}
\def\proetale{{pro\text{-}\acute{e}tale}}
\def\etale{{\acute{e}tale}}
\def\QCoh{\mathit{QCoh}}
\def\Ker{\mathop{\mathrm{Ker}}}
\def\Im{\mathop{\mathrm{Im}}}
\def\Coker{\mathop{\mathrm{Coker}}}
\def\Coim{\mathop{\mathrm{Coim}}}

% Boxtimes
%
\DeclareMathSymbol{\boxtimes}{\mathbin}{AMSa}{"02}

%
% Macros for moduli stacks/spaces
%
\def\QCohstack{\mathcal{QC}\!\mathit{oh}}
\def\Cohstack{\mathcal{C}\!\mathit{oh}}
\def\Spacesstack{\mathcal{S}\!\mathit{paces}}
\def\Quotfunctor{\mathrm{Quot}}
\def\Hilbfunctor{\mathrm{Hilb}}
\def\Curvesstack{\mathcal{C}\!\mathit{urves}}
\def\Polarizedstack{\mathcal{P}\!\mathit{olarized}}
\def\Complexesstack{\mathcal{C}\!\mathit{omplexes}}
% \Pic is the operator that assigns to X its picard group, usage \Pic(X)
% \Picardstack_{X/B} denotes the Picard stack of X over B
% \Picardfunctor_{X/B} denotes the Picard functor of X over B
\def\Pic{\mathop{\mathrm{Pic}}\nolimits}
\def\Picardstack{\mathcal{P}\!\mathit{ic}}
\def\Picardfunctor{\mathrm{Pic}}
\def\Deformationcategory{\mathcal{D}\!\mathit{ef}}

\setlength{\emergencystretch}{3em}
\begin{document}
\title[Stacks Project, Tag 0F1I]{290 --- Universal punctual base change for etale cohomology over a field}
\maketitle
\noindent Copyright (C) 2005--2025 Johan de Jong. Stacks Project authors. Source: \href{https://stacks.math.columbia.edu/tag/0F1I}{Tag 0F1I}.

\noindent Editorial difficulty note (not author text). Difficulty H2: The result treats arbitrary schemes over the field and arbitrary quasi-compact quasi-separated test morphisms. The proof approximates the scheme and coefficients, reduces hypothetical failures to algebraically closed fraction fields, then builds a directed system of finite-type opens to recover base change from the open-immersion product theorem..

\noindent Editorial provenance note (not author text). History: excerpt from revision \texttt{a0444\allowbreak 6e57e\allowbreak c1fbc\allowbreak 252a8\allowbreak 71afc\allowbreak ec775\allowbreak 2fb28\allowbreak 07b14}, etale-cohomology.tex, lines 20596--20703. Reference presentation was adapted; original-author context and core support blocks were retained or added. The mathematical wording is unchanged. Licensed under GNU FDL 1.2 or later; no invariant sections or cover texts. See ../licenses/COPYING. Context blocks reproduce author definitions and supporting results; they are not additional counted problems.

\begin{remark}
\label{remark-base-change-holds}
Let $f : X \to S$ be a morphism of schemes. Let $n$ be an integer.
We will say $BC(f, n, q_0)$ is true if for every commutative diagram
$$
\xymatrix{
X \ar[d]_f & X' \ar[l] \ar[d]_{f'} & Y \ar[l]^h \ar[d]^e \\
S & S' \ar[l] & T \ar[l]_g
}
$$
with $X' = X \times_S S'$ and $Y = X' \times_{S'} T$ and
$g$ quasi-compact and quasi-separated, and every abelian sheaf
$\mathcal{F}$ on $T_\etale$ annihilated by $n$ the base change map
$$
(f')^{-1}R^qg_*\mathcal{F}
\longrightarrow
R^qh_*e^{-1}\mathcal{F}
$$
is an isomorphism for $q \leq q_0$.
\end{remark}

\begin{definition}
\label{definition-inverse-system-sheaves}
Let $I$ be a preordered set. Let $(X_i, f_{i'i})$ be an inverse
system of schemes over $I$.
A {\it system $(\mathcal{F}_i, \varphi_{i'i})$ of sheaves
on $(X_i, f_{i'i})$} is given by
\begin{enumerate}
\item a sheaf $\mathcal{F}_i$ on $(X_i)_\etale$ for all $i \in I$,
\item for $i' \geq i$ a map
$\varphi_{i'i} : f_{i'i}^{-1}\mathcal{F}_i \to \mathcal{F}_{i'}$
of sheaves on $(X_{i'})_\etale$
\end{enumerate}
such that $\varphi_{i''i} = \varphi_{i''i'} \circ f_{i'' i'}^{-1}\varphi_{i'i}$
whenever $i'' \geq i' \geq i$.
\end{definition}

\begin{definition}
\label{algebra-definition-perfect}
Let $k$ be a field. We say $k$ is {\it perfect}
if every field extension of $k$ is separable over $k$.
\end{definition}

\begin{definition}
\label{algebra-definition-perfection}
Let $k$ be a field. The field extension $k'/k$ of Lemma \href{https://stacks.math.columbia.edu/tag/046W}{lemma perfection (Tag 046W)}
is called the {\it perfect closure} of $k$. Notation $k^{perf}/k$.
\end{definition}

\begin{lemma}
\label{lemma-kunneth-localize-on-X}
Let $K$ be a field. Let $j : U \to X$ be an open immersion of
schemes of finite type over $K$. Let $Y$ be a scheme of finite type
over $K$. Consider the diagram
$$
\xymatrix{
Y \times_{\Spec(K)} X \ar[d]_q  &
Y \times_{\Spec(K)} U \ar[l]^h \ar[d]^p \\
X & U \ar[l]_j
}
$$
Then the base change map $q^{-1}Rj_*\mathcal{F} \to Rh_*p^{-1}\mathcal{F}$
is an isomorphism for $\mathcal{F}$ an abelian sheaf on $U_\etale$
whose stalks are torsion of orders invertible in $K$.
\end{lemma}

\begin{proof}
Write $\mathcal{F} = \colim \mathcal{F}[n]$ where the colimit
is over the multiplicative system of integers invertible in $K$.
Since cohomology commutes with filtered colimits in our situation
(for a precise reference see Lemma \ref{lemma-base-change-Rf-star-colim}),
it suffices to prove the lemma for $\mathcal{F}[n]$. Thus we may
assume $\mathcal{F}$ is a sheaf of $\mathbf{Z}/n\mathbf{Z}$-modules
for some $n$ invertible in $K$ (we will use this at the very end of
the proof). In the proof we use the short hand $X \times_K Y$ for the fibre
product over $\Spec(K)$.
We will prove the lemma by induction on $\dim(X) + \dim(Y)$.
The lemma is trivial if $\dim(X) \leq 0$, since in this case
$U$ is an open and closed subscheme of $X$.
Choose a point $z \in X \times_K Y$. We will show
the stalk at $\overline{z}$ is an isomorphism.

\medskip\noindent
Suppose that $z \mapsto x \in X$ and assume $\text{trdeg}_K(\kappa(x)) > 0$.
Set $X' = \Spec(\mathcal{O}_{X, x}^{sh})$
and denote $U' \subset X'$ the inverse image of $U$.
Consider the base change
$$
\xymatrix{
Y \times_K X' \ar[d]_{q'}  &
Y \times_K U' \ar[l]^{h'} \ar[d]^{p'} \\
X' & U' \ar[l]_{j'}
}
$$
of our diagram by $X' \to X$. Observe that $X' \to X$ is a filtered
colimit of \'etale morphisms. By smooth base change in the form of
Lemma \href{https://stacks.math.columbia.edu/tag/0F09}{lemma smooth base change general (Tag 0F09)} the pullback of
$q^{-1}Rj_*\mathcal{F} \to Rh_*p^{-1}\mathcal{F}$ to $X'$ to
$Y \times_K X'$ is the map
$(q')^{-1}Rj'_*\mathcal{F}' \to Rj'_*(p')^{-1}\mathcal{F}'$ where
$\mathcal{F}'$ is the pullback of $\mathcal{F}$ to $U'$.
(In this step it would suffice to use \'etale base change which is
an essentially trivial result.) So it suffices to show
that $(q')^{-1}Rj'_*\mathcal{F}' \to Rj'_*(p')^{-1}\mathcal{F}'$
is an isomorphism in order to prove that our original
map is an isomorphism on stalks at $\overline{z}$.
By Lemma \ref{lemma-strictly-henselian}
there is a separably closed field $L/K$ such that
$X' = \lim X_i$ with $X_i$ affine of finite type over $L$
and $\dim(X_i) < \dim(X)$. For $i$ large enough there
exists an open $U_i \subset X_i$ restricting to $U'$ in $X'$.
We may apply the induction hypothesis to the diagram
$$
\vcenter{
\xymatrix{
Y \times_K X_i \ar[d]_{q_i}  &
Y \times_K U_i \ar[l]^{h_i} \ar[d]^{p_i} \\
X_i & U_i \ar[l]_{j_i}
}
}
\quad\text{equal to}\quad
\vcenter{
\xymatrix{
Y_L \times_L X_i \ar[d]_{q_i}  &
Y_L \times_L U_i \ar[l]^{h_i} \ar[d]^{p_i} \\
X_i & U_i \ar[l]_{j_i}
}
}
$$
over the field $L$ and the pullback of $\mathcal{F}$ to these diagrams.
By Lemma \ref{lemma-base-change-Rf-star-colim} we conclude that the map
$(q')^{-1}Rj'_*\mathcal{F}' \to Rj'_*(p')^{-1}\mathcal{F}$ is an isomorphism.

\medskip\noindent
Suppose that $z \mapsto y \in Y$ and assume $\text{trdeg}_K(\kappa(y)) > 0$.
Let $Y' = \Spec(\mathcal{O}_{X, x}^{sh})$.
By Lemma \ref{lemma-strictly-henselian} there is a separably closed field
$L/K$ such that $Y' = \lim Y_i$ with $Y_i$ affine of finite type over $L$
and $\dim(Y_i) < \dim(Y)$. In particular $Y'$ is a scheme over $L$.
Denote with a subscript $L$ the base change from schemes over $K$
to schemes over $L$. Consider the commutative diagrams
$$
\vcenter{
\xymatrix{
Y' \times_K X \ar[d]_f  &
Y' \times_K U \ar[l]^{h'} \ar[d]^{f'} \\
Y \times_K X \ar[d]_q  &
Y \times_K U \ar[l]^h \ar[d]^p \\
X & U \ar[l]_j
}
}
\quad\text{and}\quad
\vcenter{
\xymatrix{
Y' \times_L X_L \ar[d]_{q'}  &
Y' \times_L U_L \ar[l]^{h'} \ar[d]^{p'} \\
X_L \ar[d] &
U_L \ar[l]^{j_L} \ar[d] \\
X & U \ar[l]_j
}
}
$$
and observe the top and bottom rows are the same on the left and the right.
By smooth base change we see that
$f^{-1}Rh_*p^{-1}\mathcal{F} = Rh'_*(f')^{-1}p^{-1}\mathcal{F}$
(similarly to the previous paragraph).
By smooth base change for $\Spec(L) \to \Spec(K)$
(Lemma \ref{lemma-base-change-field-extension})
we see that $Rj_{L, *}\mathcal{F}_L$ is the pullback of
$Rj_*\mathcal{F}$ to $X_L$. Combining these two observations,
we conclude that it suffices to prove the base change map
for the upper square in the diagram on the right
is an isomorphism in order to prove that our original
map is an isomorphism on stalks at $\overline{z}$\footnote{Here we
use that a ``vertical composition'' of base change maps is a base change
map as explained in Cohomology on Sites, Remark
\href{https://stacks.math.columbia.edu/tag/0E46}{remark compose base change (Tag 0E46)}.}.
Then using that $Y' = \lim Y_i$ and argueing exactly
as in the previous paragraph we see that the induction
hypothesis forces our map over $Y' \times_K X$ to be
an isomorphism.

\medskip\noindent
Thus any counter example with $\dim(X) + \dim(Y)$ minimal
would only have nonisomorphisms
$q^{-1}Rj_*\mathcal{F} \to Rh_*p^{-1}\mathcal{F}$
on stalks at closed points of $X \times_K Y$ (because a point
$z$ of $X \times_K Y$ is a closed point if and only if
both the image of $z$ in $X$ and in $Y$ are closed).
Since it is enough to prove the isomorphism locally,
we may assume $X$ and $Y$ are affine. However, then
we can choose an open dense immersion $Y \to Y'$ with $Y'$
projective. (Choose a closed immersion $Y \to \mathbf{A}^n_K$
and let $Y'$ be the scheme theoretic closure of $Y$ in
$\mathbf{P}^n_K$.) Then $\dim(Y') = \dim(Y)$ and hence
we get a ``minimal'' counter example with $Y$ projective over $K$.
In the next paragraph we show that this can't happen.

\medskip\noindent
Consider a diagram as in the statement of the lemma
such that $q^{-1}Rj_*\mathcal{F} \to Rh_*p^{-1}\mathcal{F}$
is an isomorphism at all non-closed points of $X \times_K Y$
and such that $Y$ is projective.
The restriction of the map to $(X \times_K Y)_{K^{sep}}$
is the corresponding map for the diagram of the lemma
base changed to $K^{sep}$. Thus we may and do assume $K$
is separably algebraically closed. Choose a distinguished triangle
$$
q^{-1}Rj_*\mathcal{F} \to Rh_*p^{-1}\mathcal{F} \to Q \to
(q^{-1}Rj_*\mathcal{F})[1]
$$
in $D((X \times_K Y)_\etale)$. Since $Q$ is supported in
closed points we see that it suffices to prove
$H^i(X \times_K Y, Q) = 0$ for all $i$, see
Lemma \ref{lemma-vanishing-closed-points}.
Thus it suffices to prove that
$q^{-1}Rj_*\mathcal{F} \to Rh_*p^{-1}\mathcal{F}$
induces an isomorphism on cohomology. Recall that $\mathcal{F}$
is annihilated by $n$ invertible in $K$.
By the K\"unneth formula of Lemma \ref{lemma-kunneth-one-proper}
we have
\begin{align*}
R\Gamma(X \times_K Y, q^{-1}Rj_*\mathcal{F})
&=
R\Gamma(X, Rj_*\mathcal{F})
\otimes_{\mathbf{Z}/n\mathbf{Z}}^\mathbf{L}
R\Gamma(Y, \mathbf{Z}/n\mathbf{Z}) \\
& =
R\Gamma(U, \mathcal{F})
\otimes_{\mathbf{Z}/n\mathbf{Z}}^\mathbf{L}
R\Gamma(Y, \mathbf{Z}/n\mathbf{Z})
\end{align*}
and
$$
R\Gamma(X \times_K Y, Rh_*p^{-1}\mathcal{F}) =
R\Gamma(U \times_K Y, p^{-1}\mathcal{F}) =
R\Gamma(U, \mathcal{F})
\otimes_{\mathbf{Z}/n\mathbf{Z}}^\mathbf{L}
R\Gamma(Y, \mathbf{Z}/n\mathbf{Z})
$$
This finishes the proof.
\end{proof}


\begin{lemma}
\label{lemma-base-change-does-not-hold-post}
With $f : X \to S$ and $n$ as in Remark \ref{remark-base-change-holds}
assume $n$ is invertible on $S$ and that for some $q \geq 1$
we have that $BC(f, n, q - 1)$ is true, but $BC(f, n, q)$ is not.
Then there exist a commutative diagram
$$
\xymatrix{
X \ar[d]_f & X' \ar[d] \ar[l] & Y \ar[l]^h \ar[d] \\
S & S' \ar[l] & \Spec(K) \ar[l]
}
$$
with both squares cartesian, where $S'$ is affine, integral, and normal
with algebraically closed function field $K$ and there exists an integer
$d | n$ such that $R^qh_*(\mathbf{Z}/d\mathbf{Z})$ is nonzero.
\end{lemma}

\begin{proof}
First choose a diagram and $\mathcal{F}$ as in
Lemma \ref{lemma-base-change-does-not-hold}.
We may and do assume $S'$ is affine (this is obvious, but
see proof of the lemma in case of doubt).
Let $K'$ be the function field of $S'$
and let $Y' = X' \times_{S'} \Spec(K')$ to get the diagram
$$
\xymatrix{
X \ar[d]_f &
X' \ar[d] \ar[l] &
Y' \ar[l]^{h'} \ar[d] &
Y \ar[l] \ar[d] \\
S &
S' \ar[l] &
\Spec(K') \ar[l] &
\Spec(K) \ar[l]
}
$$
By Lemma \ref{lemma-smooth-base-change-separably-closed}
the total direct image $R(Y \to Y')_*\mathbf{Z}/d\mathbf{Z}$
is isomorphic to $\mathbf{Z}/d\mathbf{Z}$ in $D(Y'_\etale)$; here
we use that $n$ is invertible on $S$.
Thus $Rh'_*\mathbf{Z}/d\mathbf{Z} = Rh_*\mathbf{Z}/d\mathbf{Z}$
by the relative Leray spectral sequence. This finishes the proof.
\end{proof}


\begin{lemma}
\label{lemma-base-change-q-injective}
With $f : X \to S$ and $n$ as in Remark \ref{remark-base-change-holds}
assume for some $q \geq 1$ we have $BC(f, n, q - 1)$. Then
for every commutative diagram
$$
\xymatrix{
X \ar[d]_f & X' \ar[l] \ar[d]_{f'} & Y \ar[l]^h \ar[d]^e \\
S & S' \ar[l] & T \ar[l]_g
}
$$
with $X' = X \times_S S'$ and $Y = X' \times_{S'} T$ and
$g$ quasi-compact and quasi-separated, and every abelian sheaf
$\mathcal{F}$ on $T_\etale$ annihilated by $n$
\begin{enumerate}
\item the base change map
$(f')^{-1}R^qg_*\mathcal{F}\to R^qh_*e^{-1}\mathcal{F}$
is injective,
\item if $\mathcal{F} \subset \mathcal{G}$ where $\mathcal{G}$
on $T_\etale$ is annihilated by $n$, then
$$
\Coker\left(
(f')^{-1}R^qg_*\mathcal{F}\to R^qh_*e^{-1}\mathcal{F}
\right)
\subset
\Coker\left(
(f')^{-1}R^qg_*\mathcal{G}\to R^qh_*e^{-1}\mathcal{G}
\right)
$$
\item if in (2) the sheaf $\mathcal{G}$ is an injective sheaf
of $\mathbf{Z}/n\mathbf{Z}$-modules, then
$$
\Coker\left((f')^{-1}R^qg_*\mathcal{F}\to R^qh_*e^{-1}\mathcal{F} \right)
\subset R^qh_*e^{-1}\mathcal{G}
$$
\end{enumerate}
\end{lemma}

\begin{proof}
Choose a short exact sequence
$0 \to \mathcal{F} \to \mathcal{I} \to \mathcal{Q} \to 0$
where $\mathcal{I}$ is an injective sheaf of $\mathbf{Z}/n\mathbf{Z}$-modules.
Consider the induced diagram
$$
\xymatrix{
(f')^{-1}R^{q - 1}g_*\mathcal{I} \ar[d]_{\cong} \ar[r] &
(f')^{-1}R^{q - 1}g_*\mathcal{Q} \ar[d]_{\cong} \ar[r] &
(f')^{-1}R^qg_*\mathcal{F} \ar[d] \ar[r] &
0 \ar[d] \\
R^{q - 1}h_*e^{-1}\mathcal{I} \ar[r] &
R^{q - 1}h_*e^{-1}\mathcal{Q} \ar[r] &
R^qh_*e^{-1}\mathcal{F} \ar[r] &
R^qh_*e^{-1}\mathcal{I}
}
$$
with exact rows. We have the zero in the right upper corner
as $\mathcal{I}$ is injective. The left two vertical arrows are
isomorphisms by $BC(f, n, q - 1)$. We conclude that part (1) holds.
The above also shows that
$$
\Coker\left(
(f')^{-1}R^qg_*\mathcal{F}\to R^qh_*e^{-1}\mathcal{F}
\right)
\subset
R^qh_*e^{-1}\mathcal{I}
$$
hence part (3) holds. To prove (2) choose
$\mathcal{F} \subset \mathcal{G} \subset \mathcal{I}$.
\end{proof}


\begin{lemma}
\label{lemma-base-change-q-integral-top}
With $f : X \to S$ and $n$ as in Remark \ref{remark-base-change-holds}
assume for some $q \geq 1$ we have $BC(f, n, q - 1)$. Consider
commutative diagrams
$$
\vcenter{
\xymatrix{
X \ar[d]_f &
X' \ar[d]_{f'} \ar[l] &
Y \ar[l]^h \ar[d]^e &
Y' \ar[l]^{\pi'} \ar[d]^{e'} \\
S &
S' \ar[l] &
T \ar[l]_g &
T' \ar[l]_\pi
}
}
\quad\text{and}\quad
\vcenter{
\xymatrix{
X' \ar[d]_{f'} & & Y' \ar[ll]^{h' = h \circ \pi'} \ar[d]^{e'} \\
S' & & T' \ar[ll]_{g' = g \circ \pi}
}
}
$$
where all squares are cartesian, $g$ quasi-compact and quasi-separated, and
$\pi$ is integral surjective. Let $\mathcal{F}$ be an abelian sheaf
on $T_\etale$ annihilated by $n$ and set $\mathcal{F}' = \pi^{-1}\mathcal{F}$.
If the base change map
$$
(f')^{-1}R^qg'_*\mathcal{F}' \longrightarrow R^qh'_*(e')^{-1}\mathcal{F}'
$$
is an isomorphism, then the base change map
$(f')^{-1}R^qg_*\mathcal{F} \to R^qh_*e^{-1}\mathcal{F}$
is an isomorphism.
\end{lemma}

\begin{proof}
Observe that $\mathcal{F} \to \pi_*\pi^{-1}\mathcal{F}'$ is injective
as $\pi$ is surjective (check on stalks). Thus by
Lemma \ref{lemma-base-change-q-injective}
we see that it suffices to show that the base change map
$$
(f')^{-1}R^qg_*\pi_*\mathcal{F}'
\longrightarrow
R^qh_*e^{-1}\pi_*\mathcal{F}'
$$
is an isomorphism. This follows from the assumption because
we have $R^qg_*\pi_*\mathcal{F}' = R^qg'_*\mathcal{F}'$,
we have $e^{-1}\pi_*\mathcal{F}' =\pi'_*(e')^{-1}\mathcal{F}'$, and
we have $R^qh_*\pi'_*(e')^{-1}\mathcal{F}' = R^qh'_*(e')^{-1}\mathcal{F}'$.
This follows from Lemmas
\href{https://stacks.math.columbia.edu/tag/0EYP}{lemma integral pushforward commutes with base change (Tag 0EYP)} and
\href{https://stacks.math.columbia.edu/tag/04C2}{lemma what integral (Tag 04C2)} and the relative leray spectral sequence
(Cohomology on Sites, Lemma \href{https://stacks.math.columbia.edu/tag/0734}{lemma relative Leray (Tag 0734)}).
\end{proof}


\begin{lemma}
\label{lemma-base-change-q-integral-bottom}
With $f : X \to S$ and $n$ as in Remark \ref{remark-base-change-holds}
assume for some $q \geq 1$ we have $BC(f, n, q - 1)$. Consider
commutative diagrams
$$
\vcenter{
\xymatrix{
X \ar[d]_f &
X' \ar[d]_{f'} \ar[l] &
X'' \ar[l]^{\pi'} \ar[d]_{f''} &
Y \ar[l]^{h'} \ar[d]^e \\
S &
S' \ar[l] &
S'' \ar[l]_\pi &
T \ar[l]_{g'}
}
}
\quad\text{and}\quad
\vcenter{
\xymatrix{
X' \ar[d]_{f'} & & Y \ar[ll]^{h = h' \circ \pi'} \ar[d]^e \\
S' & & T \ar[ll]_{g = g' \circ \pi}
}
}
$$
where all squares are cartesian, $g'$ quasi-compact and quasi-separated, and
$\pi$ is integral. Let $\mathcal{F}$ be an abelian sheaf
on $T_\etale$ annihilated by $n$. If the base change map
$$
(f')^{-1}R^qg_*\mathcal{F} \longrightarrow R^qh_*e^{-1}\mathcal{F}
$$
is an isomorphism, then the base change map
$(f'')^{-1}R^qg'_*\mathcal{F} \to R^qh'_*e^{-1}\mathcal{F}$
is an isomorphism.
\end{lemma}

\begin{proof}
Since $\pi$ and $\pi'$ are integral we have $R\pi_* = \pi_*$ and
$R\pi'_* = \pi'_*$, see Lemma \href{https://stacks.math.columbia.edu/tag/04C2}{lemma what integral (Tag 04C2)}.
We also have $(f')^{-1}\pi_* = \pi'_*(f'')^{-1}$. Thus we see that
$\pi'_*(f'')^{-1}R^qg'_*\mathcal{F} = (f')^{-1}R^qg_*\mathcal{F}$
and
$\pi'_*R^qh'_*e^{-1}\mathcal{F} = R^qh_*e^{-1}\mathcal{F}$.
Thus the assumption means that our map becomes an
isomorphism after applying the functor $\pi'_*$.
Hence we see that it is an isomorphism by Lemma \href{https://stacks.math.columbia.edu/tag/04C2}{lemma what integral (Tag 04C2)}.
\end{proof}


\begin{lemma}
\label{lemma-formal-argument}
Let $T$ be a quasi-compact and quasi-separated scheme.
Let $P$ be a property for quasi-compact and quasi-separated
schemes over $T$. Assume
\begin{enumerate}
\item If $T'' \to T'$ is a thickening of quasi-compact and
quasi-separated schemes over $T$, then $P(T'')$ if and only if $P(T')$.
\item If $T' = \lim T_i$ is a limit of an inverse system of
quasi-compact and quasi-separated schemes over $T$ with affine
transition morphisms and $P(T_i)$ holds for all $i$, then
$P(T')$ holds.
\item If $Z \subset T'$ is a closed subscheme with
quasi-compact complement $V \subset T'$ and $P(T')$ holds,
then either $P(V)$ or $P(Z)$ holds.
\end{enumerate}
Then $P(T)$ implies $P(\Spec(K))$ for some morphism $\Spec(K) \to T$
where $K$ is a field.
\end{lemma}

\begin{proof}
Consider the set $\mathfrak T$ of closed subschemes $T' \subset T$
such that $P(T')$. By assumption (2) this set has a minimal element,
say $T'$. By assumption (1) we see that $T'$ is reduced.
Let $\eta \in T'$ be the generic point of an irreducible
component of $T'$. Then $\eta = \Spec(K)$ for some field
$K$ and $\eta = \lim V$ where the limit is over the affine
open subschemes $V \subset T'$ containing $\eta$.
By assumption (3) and the minimality of $T'$ we see
that $P(V)$ holds for all these $V$. Hence $P(\eta)$
by (2) and the proof is complete.
\end{proof}


\begin{lemma}
\label{lemma-base-change-does-not-hold-pre}
With $f : X \to S$ and $n$ as in Remark \ref{remark-base-change-holds}
assume for some $q \geq 1$ we have that $BC(f, n, q - 1)$ is true, but
$BC(f, n, q)$ is not. Then there exist a commutative diagram
$$
\xymatrix{
X \ar[d]_f & X' \ar[d]_{f'} \ar[l] & Y \ar[l]^h \ar[d]^e \\
S & S' \ar[l] & \Spec(K) \ar[l]_g
}
$$
where $X' = X \times_S S'$, $Y = X' \times_{S'} \Spec(K)$,
$K$ is a field, and $\mathcal{F}$ is an abelian sheaf
on $\Spec(K)$ annihilated by $n$ such that
$(f')^{-1}R^qg_*\mathcal{F} \to R^qh_*e^{-1}\mathcal{F}$
is not an isomorphism.
\end{lemma}

\begin{proof}
Choose a commutative diagram
$$
\xymatrix{
X \ar[d]_f & X' \ar[l] \ar[d]_{f'} & Y \ar[l]^h \ar[d]^e \\
S & S' \ar[l] & T \ar[l]_g
}
$$
with $X' = X \times_S S'$ and $Y = X' \times_{S'} T$ and
$g$ quasi-compact and quasi-separated, and an abelian sheaf
$\mathcal{F}$ on $T_\etale$ annihilated by $n$ such that
the base change map
$(f')^{-1}R^qg_*\mathcal{F} \to R^qh_*e^{-1}\mathcal{F}$
is not an isomorphism. Of course we may and do replace $S'$
by an affine open of $S'$; this implies that $T$ is quasi-compact
and quasi-separated. By Lemma \ref{lemma-base-change-q-injective} we see
$(f')^{-1}R^qg_*\mathcal{F} \to R^qh_*e^{-1}\mathcal{F}$
is injective. Pick a geometric point $\overline{x}$ of $X'$
and an element $\xi$ of $(R^qh_*q^{-1}\mathcal{F})_{\overline{x}}$
which is not in the image of the map
$((f')^{-1}R^qg_*\mathcal{F})_{\overline{x}} \to
(R^qh_*e^{-1}\mathcal{F})_{\overline{x}}$.

\medskip\noindent
Consider a morphism $\pi : T' \to T$ with $T'$ quasi-compact
and quasi-separated and denote $\mathcal{F}' = \pi^{-1}\mathcal{F}$.
Denote $\pi' : Y' = Y \times_T T' \to Y$ the base change of $\pi$
and $e' : Y' \to T'$ the base change of $e$. Picture
$$
\vcenter{
\xymatrix{
X' \ar[d]_{f'} & Y \ar[l]^h \ar[d]^e & Y' \ar[l]^{\pi'} \ar[d]^{e'} \\
S' & T \ar[l]_g & T' \ar[l]_\pi
}
}
\quad\text{and}\quad
\vcenter{
\xymatrix{
X' \ar[d]_{f'} & & Y' \ar[ll]^{h' = h \circ \pi'} \ar[d]^{e'} \\
S' & & T' \ar[ll]_{g' = g \circ \pi}
}
}
$$
Using pullback maps we obtain a canonical commutative diagram
$$
\xymatrix{
(f')^{-1}R^qg_*\mathcal{F} \ar[r] \ar[d] &
(f')^{-1}R^qg'_*\mathcal{F}' \ar[d] \\
R^qh_*e^{-1}\mathcal{F} \ar[r] &
R^qh'_*(e')^{-1}\mathcal{F}'
}
$$
of abelian sheaves on $X'$. Let $P(T')$ be the property
\begin{itemize}
\item The image $\xi'$ of $\xi$ in
$(Rh'_*(e')^{-1}\mathcal{F}')_{\overline{x}}$ is
not in the image of the map
$(f^{-1}R^qg'_*\mathcal{F}')_{\overline{x}} \to
(R^qh'_*(e')^{-1}\mathcal{F}')_{\overline{x}}$.
\end{itemize}
We claim that hypotheses (1), (2), and (3) of
Lemma \ref{lemma-formal-argument} hold for $P$
which proves our lemma.

\medskip\noindent
Condition (1) of Lemma \ref{lemma-formal-argument}
holds for $P$ because the \'etale topology of a scheme
and a thickening of the scheme is the same. See
Proposition \href{https://stacks.math.columbia.edu/tag/03SI}{proposition topological invariance (Tag 03SI)}.

\medskip\noindent
Suppose that $I$ is a directed set and that $T_i$
is an inverse system over $I$ of quasi-compact and quasi-separated
schemes over $T$ with affine transition morphisms.
Set $T' = \lim T_i$. Denote $\mathcal{F}'$ and $\mathcal{F}_i$
the pullback of $\mathcal{F}$ to $T'$, resp.\ $T_i$. Consider
the diagrams
$$
\vcenter{
\xymatrix{
X \ar[d]_{f'} & Y \ar[l]^h \ar[d]^e & Y_i \ar[l]^{\pi_i'} \ar[d]^{e_i} \\
S & T \ar[l]_g & T_i \ar[l]_{\pi_i}
}
}
\quad\text{and}\quad
\vcenter{
\xymatrix{
X \ar[d]_{f'} & & Y_i \ar[ll]^{h_i = h \circ \pi_i'} \ar[d]^{e_i} \\
S & & T_i \ar[ll]_{g_i = g \circ \pi_i}
}
}
$$
as in the previous paragraph. It is clear that $\mathcal{F}'$ on
$T'$ is the colimit of the pullbacks of $\mathcal{F}_i$ to $T'$
and that $(e')^{-1}\mathcal{F}'$ is the colimit of the pullbacks
of $e_i^{-1}\mathcal{F}_i$ to $Y'$.
By Lemma \ref{lemma-relative-colimit-general}
we have
$$
R^qh'_*(e')^{-1}\mathcal{F}' = \colim R^qh_{i, *}e_i^{-1}\mathcal{F}_i
\quad\text{and}\quad
(f')^{-1}R^qg'_*\mathcal{F}' = \colim (f')^{-1}R^qg_{i, *}\mathcal{F}_i
$$
It follows that if $P(T_i)$ is true for all $i$, then
$P(T')$ holds. Thus condition (2) of Lemma \ref{lemma-formal-argument}
holds for $P$.

\medskip\noindent
The most interesting is condition (3) of Lemma \ref{lemma-formal-argument}.
Assume $T'$ is a quasi-compact and quasi-separated scheme over $T$
such that $P(T')$ is true.
Let $Z \subset T'$ be a closed subscheme with complement $V \subset T'$
quasi-compact. Consider the diagram
$$
\xymatrix{
Y' \times_{T'} Z \ar[d]_{e_Z} \ar[r]_{i'} &
Y' \ar[d]_{e'} &
Y' \times_{T'} V \ar[l]^{j'} \ar[d]^{e_V} \\
Z \ar[r]^i &
T' &
V \ar[l]_j
}
$$
Choose an injective map $j^{-1}\mathcal{F}' \to \mathcal{J}$
where $\mathcal{J}$ is an injective sheaf of $\mathbf{Z}/n\mathbf{Z}$-modules
on $V$. Looking at stalks we see that the map
$$
\mathcal{F}' \to \mathcal{G} = j_*\mathcal{J} \oplus i_*i^{-1}\mathcal{F}'
$$
is injective. Thus $\xi'$ maps to a nonzero element of
\begin{align*}
& \Coker\left(
((f')^{-1}R^qg'_*\mathcal{G})_{\overline{x}}
\to
(R^qh'_*(e')^{-1}\mathcal{G})_{\overline{x}}
\right) = \\
&
\Coker\left(
((f')^{-1}R^qg'_*j_*\mathcal{J})_{\overline{x}}
\to
(R^qh'_*(e')^{-1}j_*\mathcal{J})_{\overline{x}}
\right) \oplus \\
& \Coker\left(
((f')^{-1}R^qg'_*i_*i^{-1}\mathcal{F}')_{\overline{x}}
\to
(R^qh'_*(e')^{-1}i_*i^{-1}\mathcal{F}')_{\overline{x}}
\right)
\end{align*}
by part (2) of Lemma \ref{lemma-base-change-q-injective}.
If $\xi'$ does not map to zero in the second summand, then
we use
$$
(f')^{-1}R^qg'_*i_*i^{-1}\mathcal{F}' =
(f')^{-1}R^q(g' \circ i)_*i^{-1}\mathcal{F}'
$$
(because $Ri_* = i_*$ by
Proposition \href{https://stacks.math.columbia.edu/tag/03QP}{proposition finite higher direct image zero (Tag 03QP)}) and
$$
R^qh'_*(e')^{-1}i_*i^{-1}\mathcal{F} =
R^qh'_*i'_*e_Z^{-1}i^{-1}\mathcal{F} =
R^q(h' \circ i')_*e_Z^{-1}i^{-1}\mathcal{F}'
$$
(first equality by
Lemma \href{https://stacks.math.columbia.edu/tag/0959}{lemma finite pushforward commutes with base change (Tag 0959)}
and the second because
$Ri'_* = i'_*$ by
Proposition \href{https://stacks.math.columbia.edu/tag/03QP}{proposition finite higher direct image zero (Tag 03QP)})
to we see that we have $P(Z)$.
Finally, suppose $\xi'$ does not map to zero in the first summand.
We have
$$
(e')^{-1}j_*\mathcal{J} = j'_*e_V^{-1}\mathcal{J}
\quad\text{and}\quad
R^aj'_*e_V^{-1}\mathcal{J} = 0, \quad a = 1, \ldots, q - 1
$$
by $BC(f, n, q - 1)$ applied to the diagram
$$
\xymatrix{
X \ar[d]_f & Y' \ar[l] \ar[d]_{e'} & Y \ar[l]^{j'} \ar[d]^{e_V} \\
S & T' \ar[l] & V \ar[l]_j
}
$$
and the fact that $\mathcal{J}$ is injective.
By the relative Leray spectral sequence for $h' \circ j'$
(Cohomology on Sites, Lemma \href{https://stacks.math.columbia.edu/tag/0734}{lemma relative Leray (Tag 0734)})
we deduce that
$$
R^qh'_*(e')^{-1}j_*\mathcal{J} =
R^qh'_*j'_*e_V^{-1}\mathcal{J}
\longrightarrow
R^q(h' \circ j')_* e_V^{-1}\mathcal{J}
$$
is injective. Thus $\xi$ maps to a nonzero element of
$(R^q(h' \circ j')_* e_V^{-1}\mathcal{J})_{\overline{x}}$.
Applying part (3) of Lemma \ref{lemma-base-change-q-injective}
to the injection $j^{-1}\mathcal{F}' \to \mathcal{J}$
we conclude that $P(V)$ holds.
\end{proof}


\begin{lemma}
\label{lemma-strictly-henselian}
Let $K$ be a field. Let $X$ be a scheme of finite type over $K$.
Let $x \in X$. Set $a = \text{trdeg}_K(\kappa(x))$
and $d = \dim_x(X)$. Then there is a map
$$
K(t_1, \ldots, t_a)^{sep} \longrightarrow \mathcal{O}_{X, x}^{sh}
$$
such that
\begin{enumerate}
\item the residue field of $\mathcal{O}_{X, x}^{sh}$ is a purely inseparable
extension of $K(t_1, \ldots, t_a)^{sep}$,
\item $\mathcal{O}_{X, x}^{sh}$ is a filtered colimit of finite
type $K(t_1, \ldots, t_a)^{sep}$-algebras of dimension $\leq d - a$.
\end{enumerate}
\end{lemma}

\begin{proof}
We may assume $X$ is affine. By Noether normalization, after possibly
shrinking $X$ again, we can choose a finite morphism
$\pi : X \to \mathbf{A}^d_K$, see
Algebra, Lemma \href{https://stacks.math.columbia.edu/tag/00OZ}{lemma Noether normalization at point (Tag 00OZ)}.
Since $\kappa(x)$ is a finite extension of the residue field of $\pi(x)$,
this residue field has transcendence degree $a$ over $K$ as well.
Thus we can find a finite morphism $\pi' : \mathbf{A}^d_K \to \mathbf{A}^d_K$
such that $\pi'(\pi(x))$ corresponds to the generic point of the linear
subspace $\mathbf{A}^a_K \subset \mathbf{A}^d_K$ given by setting
the last $d - a$ coordinates equal to zero. Hence the composition
$$
X \xrightarrow{\pi' \circ \pi} \mathbf{A}^d_K \xrightarrow{p} \mathbf{A}^a_K
$$
of $\pi' \circ \pi$ and the projection $p$ onto the first $a$ coordinates
maps $x$ to the generic point $\eta \in \mathbf{A}^a_K$. The induced map
$$
K(t_1, \ldots, t_a)^{sep} =
\mathcal{O}_{\mathbf{A}^a_k, \eta}^{sh}
\longrightarrow \mathcal{O}_{X, x}^{sh}
$$
on \'etale local rings satisfies (1) since it is clear that the residue field
of $\mathcal{O}_{X, x}^{sh}$ is an algebraic extension of the separably closed
field $K(t_1, \ldots, t_a)^{sep}$. On the other hand, if $X = \Spec(B)$,
then $\mathcal{O}_{X, x}^{sh} = \colim B_j$ is a filtered colimit
of \'etale $B$-algebras $B_j$. Observe that $B_j$ is quasi-finite over
$K[t_1, \ldots, t_d]$ as $B$ is finite over $K[t_1, \ldots, t_d]$.
We may similarly write
$K(t_1, \ldots, t_a)^{sep} = \colim A_i$ as a filtered colimit
of \'etale $K[t_1, \ldots, t_a]$-algebras. For every $i$ we can
find an $j$ such that
$A_i \to K(t_1, \ldots, t_a)^{sep} \to \mathcal{O}_{X, x}^{sh}$
factors through a map $\psi_{i, j} : A_i \to B_j$. Then $B_j$ is quasi-finite
over $A_i[t_{a + 1}, \ldots, t_d]$. Hence
$$
B_{i, j} = B_j \otimes_{\psi_{i, j}, A_i} K(t_1, \ldots, t_a)^{sep}
$$
has dimension $\leq d - a$ as it is quasi-finite over
$K(t_1, \ldots, t_a)^{sep}[t_{a + 1}, \ldots, t_d]$.
The proof of (2) is now finished as $\mathcal{O}_{X, x}^{sh}$ is a
filtered colimit\footnote{Let $R$ be a ring.
Let $A = \colim_{i \in I} A_i$
be a filtered colimit of finitely presented $R$-algebras.
Let $B = \colim_{j \in J} B_j$ be a filtered colimit of $R$-algebras.
Let $A \to B$ be an $R$-algebra map.
Assume that for all $i \in I$ there is a $j \in J$ and an $R$-algebra map
$\psi_{i, j} : A_i \to B_j$.
Say $(i', j', \psi_{i', j'}) \geq (i, j, \psi_{i, j})$ if
$i' \geq i$, $j' \geq j$, and $\psi_{i, j}$ and $\psi_{i', j'}$
are compatible. Then the collection of triples forms a
directed set and $B = \colim B_j \otimes_{\psi_{i, j} A_i} A$.}
of the algebras $B_{i, j}$. Some details omitted.
\end{proof}


\begin{lemma}
\label{lemma-kunneth-one-proper}
Let $k$ be a separably closed field. Let $X$ be a proper scheme over $k$.
Let $Y$ be a quasi-compact and quasi-separated scheme over $k$.
\begin{enumerate}
\item If $E \in D^+(X_\etale)$ has torsion cohomology sheaves and
$K \in D^+(Y_\etale)$, then
$$
R\Gamma(X \times_{\Spec(k)} Y,
\text{pr}_1^{-1}E
\otimes_\mathbf{Z}^\mathbf{L}
\text{pr}_2^{-1}K
)
=
R\Gamma(X, E)
\otimes_\mathbf{Z}^\mathbf{L}
R\Gamma(Y, K)
$$
\item If $n \geq 1$ is an integer, $Y$ is of finite type over $k$,
$E \in D(X_\etale, \mathbf{Z}/n\mathbf{Z})$, and
$K \in D(Y_\etale, \mathbf{Z}/n\mathbf{Z})$, then
$$
R\Gamma(X \times_{\Spec(k)} Y,
\text{pr}_1^{-1}E
\otimes_{\mathbf{Z}/n\mathbf{Z}}^\mathbf{L}
\text{pr}_2^{-1}K
)
=
R\Gamma(X, E)
\otimes_{\mathbf{Z}/n\mathbf{Z}}^\mathbf{L}
R\Gamma(Y, K)
$$
\end{enumerate}
\end{lemma}

\begin{proof}
Proof of (1). By Lemma \href{https://stacks.math.columbia.edu/tag/0F0F}{lemma projection formula proper (Tag 0F0F)} we have
$$
R\text{pr}_{2, *}(
\text{pr}_1^{-1}E
\otimes_\mathbf{Z}^\mathbf{L}
\text{pr}_2^{-1}K) =
R\text{pr}_{2, *}(\text{pr}_1^{-1}E)
\otimes_\mathbf{Z}^\mathbf{L}
K
$$
By proper base change (in the form of Lemma \href{https://stacks.math.columbia.edu/tag/0DDE}{lemma proper base change (Tag 0DDE)})
this is equal to the object
$$
\underline{R\Gamma(X, E)}
\otimes_\mathbf{Z}^\mathbf{L}
K
$$
of $D(Y_\etale)$. Taking $R\Gamma(Y, -)$ on this object reproduces the
left hand side of the equality in (1) by the Leray spectral sequence
for $\text{pr}_2$. Thus we conclude by
Lemma \href{https://stacks.math.columbia.edu/tag/0F0E}{lemma pull out constant (Tag 0F0E)}.

\medskip\noindent
Proof of (2). This is exactly the same as the proof of (1)
except that we use Lemmas \href{https://stacks.math.columbia.edu/tag/0F0G}{lemma projection formula proper mod n (Tag 0F0G)},
\href{https://stacks.math.columbia.edu/tag/0F0C}{lemma proper base change mod n (Tag 0F0C)}, and
\href{https://stacks.math.columbia.edu/tag/0F12}{lemma pull out constant mod n (Tag 0F12)} as well as
$\text{cd}(Y) < \infty$ by Lemma \href{https://stacks.math.columbia.edu/tag/0F10}{lemma finite cd (Tag 0F10)}.
\end{proof}


\begin{lemma}
\label{lemma-supported-in-closed-points}
Let $K$ be a separably closed field. Let $X$ be a scheme of finite
type over $K$. Let $\mathcal{F}$ be an abelian sheaf on $X_\etale$
whose support is contained in the set of closed points of $X$.
Then $H^q(X, \mathcal{F}) = 0$ for $q > 0$ and $\mathcal{F}$
is globally generated.
\end{lemma}

\begin{proof}
(If $\mathcal{F}$ is torsion, then the vanishing follows immediately
from Lemma \href{https://stacks.math.columbia.edu/tag/0F0W}{lemma interlude II (Tag 0F0W)}.)
By Lemma \href{https://stacks.math.columbia.edu/tag/0F0N}{lemma support in subset (Tag 0F0N)} we can write $\mathcal{F}$
as a filtered colimit of constructible sheaves $\mathcal{F}_i$
of $\mathbf{Z}$-modules whose supports $Z_i \subset X$ are finite
sets of closed points. By
Proposition \href{https://stacks.math.columbia.edu/tag/04CA}{proposition closed immersion pushforward (Tag 04CA)}
such a sheaf is of the form $(Z_i \to X)_*\mathcal{G}_i$
where $\mathcal{G}_i$ is a sheaf on $Z_i$. As $K$ is separably closed,
the scheme $Z_i$ is a finite disjoint union of spectra of separably closed
fields. Recall that $H^q(Z_i, \mathcal{G}_i) = H^q(X, \mathcal{F}_i)$
by the Leray spectral sequence for $Z_i \to X$ and vanising of higher
direct images for this morphism
(Proposition \href{https://stacks.math.columbia.edu/tag/03QP}{proposition finite higher direct image zero (Tag 03QP)}).
By Lemmas \href{https://stacks.math.columbia.edu/tag/04JQ}{lemma equivalence abelian sheaves point (Tag 04JQ)} and
\href{https://stacks.math.columbia.edu/tag/03QU}{lemma compare cohomology point (Tag 03QU)}
we see that $H^q(Z_i, \mathcal{G}_i)$ is zero for $q > 0$
and that $H^0(Z_i, \mathcal{G}_i)$ generates $\mathcal{G}_i$.
We conclude the vanishing of $H^q(X, \mathcal{F}_i)$ for $q > 0$
and that $\mathcal{F}_i$ is generated by global sections.
By Theorem \href{https://stacks.math.columbia.edu/tag/09YQ}{theorem colimit (Tag 09YQ)} we see that
$H^q(X, \mathcal{F}) = 0$ for $q > 0$.
The proof is now done because a
filtered colimit of globally generated sheaves of abelian groups
is globally generated (details omitted).
\end{proof}


\begin{lemma}
\label{lemma-vanishing-closed-points}
Let $K$ be a separably closed field. Let $X$ be a scheme of finite
type over $K$. Let $Q \in D(X_\etale)$. Assume that $Q_{\overline{x}}$
is nonzero only if $x$ is a closed point of $X$. Then
$$
Q = 0 \Leftrightarrow H^i(X, Q) = 0 \text{ for all }i
$$
\end{lemma}

\begin{proof}
The implication from left to right is trivial. Thus we need
to prove the reverse implication.

\medskip\noindent
Assume $Q$ is bounded below; this cases suffices for almost all
applications. If $Q$ is not zero, then we can look at the smallest $i$
such that the cohomology sheaf $H^i(Q)$ is nonzero.
By Lemma \ref{lemma-supported-in-closed-points} we have
$H^i(X, Q) = H^0(X, H^i(Q)) \not = 0$ and we conclude.

\medskip\noindent
General case. Let $\mathcal{B} \subset \Ob(X_\etale)$ be the
quasi-compact objects. By Lemma \ref{lemma-supported-in-closed-points}
the assumptions of Cohomology on Sites, Lemma
\href{https://stacks.math.columbia.edu/tag/0BKZ}{lemma cohomology over U trivial (Tag 0BKZ)}
are satisfied. We conclude that $H^q(U, Q) = H^0(U, H^q(Q))$
for all $U \in \mathcal{B}$. In particular, this holds for $U = X$.
Thus the conclusion by Lemma \ref{lemma-supported-in-closed-points}
as $Q$ is zero in $D(X_\etale)$ if and only if $H^q(Q)$ is zero
for all $q$.
\end{proof}


\begin{lemma}
\label{lemma-base-change-Rf-star-colim}
Let $I$ be a directed set. Consider an inverse system of
cartesian diagrams of schemes
$$
\xymatrix{
X_i \ar[d]_{f_i} & Y_i \ar[l]^{h_i} \ar[d]^{e_i} \\
S_i & T_i \ar[l]_{g_i}
}
$$
with affine transition morphisms and with $g_i$ quasi-compact and
quasi-separated. Set $X = \lim X_i$,
$S = \lim S_i$, $T = \lim T_i$ and $Y = \lim Y_i$ to
obtain the cartesian diagram
$$
\xymatrix{
X \ar[d]_f & Y \ar[l]^h \ar[d]^e \\
S & T \ar[l]_g
}
$$
Let $(\mathcal{F}_i, \varphi_{i'i})$ be a system of sheaves on
$(T_i)$ as in Definition \ref{definition-inverse-system-sheaves}. Set
$\mathcal{F} = \colim p_i^{-1}\mathcal{F}_i$ on $T$
where $p_i : T \to T_i$ is the projection.
Then we have the following
\begin{enumerate}
\item If $f_i^{-1}g_{i, *}\mathcal{F}_i = h_{i, *}e_i^{-1}\mathcal{F}_i$
for all $i$, then
$f^{-1}g_*\mathcal{F} = h_*e^{-1}\mathcal{F}$.
\item If $\mathcal{F}_i$ is an abelian sheaf for all $i$ and
$f_i^{-1}R^qg_{i, *}\mathcal{F}_i = R^qh_{i, *}e_i^{-1}\mathcal{F}_i$
for all $i$, then
$f^{-1}R^qg_*\mathcal{F} = R^qh_*e^{-1}\mathcal{F}$.
\end{enumerate}
\end{lemma}

\begin{proof}
We prove (2) and we omit the proof of (1). We will use without further
mention that pullback of sheaves commutes with colimits as it is a
left adjoint. Observe that $h_i$ is quasi-compact and quasi-separated as a
base change of $g_i$.
Denoting $q_i : Y \to Y_i$ the projections, observe that
$e^{-1}\mathcal{F} = \colim e^{-1}p_i^{-1}\mathcal{F}_i =
\colim q_i^{-1}e_i^{-1}\mathcal{F}_i$.
By Lemma \ref{lemma-relative-colimit-general}
this gives
$$
R^qh_*e^{-1}\mathcal{F} = \colim r_i^{-1}R^qh_{i, *}e_i^{-1}\mathcal{F}_i
$$
where $r_i : X \to X_i$ is the projection.
Similarly, we have
$$
f^{-1}Rg_*\mathcal{F} =
f^{-1}\colim s_i^{-1}R^qg_{i, *}\mathcal{F}_i =
\colim r_i^{-1}f_i^{-1}R^qg_{i, *}\mathcal{F}_i
$$
where $s_i : S \to S_i$ is the projection. The lemma follows.
\end{proof}


\begin{lemma}
\label{lemma-relative-colimit-general}
Let $I$ be a directed set. Let $g_i : X_i \to S_i$ be an inverse system of
morphisms of schemes over $I$. Assume $g_i$ is quasi-compact and
quasi-separated and for $i' \geq i$ the transition morphisms
$f_{i'i} : X_{i'} \to X_i$ and $h_{i'i} : S_{i'} \to S_i$ are affine.
Let $g : X \to S$ be the limit of the morphisms $g_i$, see
Limits, Section \href{https://stacks.math.columbia.edu/tag/01YV}{section limits (Tag 01YV)}.
Denote $f_i : X \to X_i$ and $h_i : S \to S_i$ the projections.
Let $(\mathcal{F}_i, \varphi_{i'i})$ be a system of sheaves
on $(X_i, f_{i'i})$. Set $\mathcal{F} = \colim f_i^{-1}\mathcal{F}_i$. Then
$$
R^p g_* \mathcal{F} =
\colim_{i \in I} h_i^{-1}R^p g_{i, *} \mathcal{F}_i
$$
for all $p \geq 0$.
\end{lemma}

\begin{proof}
How is the map of the lemma constructed?
For $i' \geq i$ we have a commutative diagram
$$
\xymatrix{
X \ar[r]_{f_{i'}} \ar[d]_g &
X_{i'} \ar[r]_{f_{i'i}} \ar[d]_{g_{i'}} &
X_i \ar[d]^{g_i} \\
S \ar[r]^{h_{i'}} &
S_{i'} \ar[r]^{h_{i'i}} &
S_i
}
$$
If we combine the base change map
$h_{i'i}^{-1}Rg_{i, *}\mathcal{F}_i \to Rg_{i', *}f_{i'i}^{-1}\mathcal{F}_i$
(Cohomology on Sites, Lemma
\href{https://stacks.math.columbia.edu/tag/0736}{lemma base change map flat case (Tag 0736)} or
Remark \href{https://stacks.math.columbia.edu/tag/07A7}{remark base change (Tag 07A7)})
with the map $Rg_{i', *}\varphi_{i'i}$, then we obtain
$\psi_{i'i} : h_{i' i}^{-1} R^p g_{i, *} \mathcal{F}_i \to
R^pg_{i', *} \mathcal{F}_{i'}$. Similarly, using the left square
in the diagram we obtain maps
$\psi_i : h_i^{-1}R^pg_{i, *}\mathcal{F}_i \to R^pg_*\mathcal{F}$.
The maps $h_{i'}^{-1}\psi_{i'i}$ and $\psi_i$ are the maps used in
the statement of the lemma. For this to make sense, we have to check that
$\psi_{i''i} = \psi_{i''i'} \circ h_{i''i'}^{-1}\psi_{i'i}$ and
$\psi_{i'} \circ h_{i'}^{-1}\psi_{i'i} = \psi_i$; this follows
from Cohomology on Sites, Remark
\href{https://stacks.math.columbia.edu/tag/0E47}{remark compose base change horizontal (Tag 0E47)}.

\medskip\noindent
Proof of the equality. First proof using
dimension shifting\footnote{You can also use this method
to produce the maps in the lemma.}. For any
$U$ affine and \'etale over $X$ by Theorem \href{https://stacks.math.columbia.edu/tag/09YQ}{theorem colimit (Tag 09YQ)}
we have
$$
g_*\mathcal{F}(U) =
H^0(U \times_S X, \mathcal{F}) =
\colim H^0(U_i \times_{S_i} X_i, \mathcal{F}_i) =
\colim g_{i, *}\mathcal{F}_i(U_i)
$$
where the colimit is over $i$ large enough such that
there exists an $i$ and $U_i$ affine \'etale over $S_i$
whose base change is $U$ over $S$ (see
Lemma \href{https://stacks.math.columbia.edu/tag/0EYL}{lemma colimit affine sites (Tag 0EYL)}).
The right hand side is equal to
$(\colim h_i^{-1}g_{i, *}\mathcal{F}_i)(U)$ by
Sites, Lemma \href{https://stacks.math.columbia.edu/tag/09YN}{lemma colimit (Tag 09YN)}.
This proves the lemma for $p = 0$.
If $(\mathcal{G}_i, \varphi_{i'i})$ is a system with
$\mathcal{G} = \colim f_i^{-1}\mathcal{G}_i$
such that $\mathcal{G}_i$ is an injective abelian sheaf on $X_i$
for all $i$, then for any $U$ affine and \'etale over $X$ by
Theorem \href{https://stacks.math.columbia.edu/tag/09YQ}{theorem colimit (Tag 09YQ)} we have
$$
H^p(U \times_S X, \mathcal{G}) =
\colim H^p(U_i \times_{S_i} X_i, \mathcal{G}_i) = 0
$$
for $p > 0$ (same colimit as before). Hence $R^pg_*\mathcal{G} = 0$
and we get the result for $p > 0$ for such a system.
In general we may choose a short exact sequence of systems
$$
0 \to (\mathcal{F}_i, \varphi_{i'i}) \to
(\mathcal{G}_i, \varphi_{i'i}) \to
(\mathcal{Q}_i, \varphi_{i'i}) \to 0
$$
where $(\mathcal{G}_i, \varphi_{i'i})$ is as above, see
Cohomology on Sites, Lemma
\href{https://stacks.math.columbia.edu/tag/0EXZ}{lemma colim sites injective (Tag 0EXZ)}.
By induction the lemma holds for $p - 1$ and by the above we have
vanishing for $p$ and $(\mathcal{G}_i, \varphi_{i'i})$.
Hence the result for $p$
and $(\mathcal{F}_i, \varphi_{i'i})$ by the long exact sequence
of cohomology.

\medskip\noindent
Second proof.
Recall that $S_{affine, \etale} = \colim (S_i)_{affine, \etale}$, see
Lemma \href{https://stacks.math.columbia.edu/tag/0EYL}{lemma colimit affine sites (Tag 0EYL)}. Thus if $U$ is an object of
$S_{affine, \etale}$, then we can write $U = U_i \times_{S_i} S$
for some $i$ and some $U_i$ in $(S_i)_{affine, \etale}$ and
$$
(\colim_{i \in I} h_i^{-1}R^p g_{i, *} \mathcal{F}_i)(U) =
\colim_{i' \geq i} (R^p g_{i', *}\mathcal{F}_{i'})(U_i \times_{S_i} S_{i'})
$$
by Sites, Lemma \href{https://stacks.math.columbia.edu/tag/09YN}{lemma colimit (Tag 09YN)} and the construction of the
transition maps in the system described above. Since
$R^pg_{i', *}\mathcal{F}_{i'}$ is the sheaf associated to the presheaf
$U_{i'} \mapsto H^p(U_{i'} \times_{S_{i'}} X_{i'}, \mathcal{F}_{i'})$
and since $R^pg_*\mathcal{F}$ is the sheaf associated to the presheaf
$U \mapsto H^p(U \times_S X, \mathcal{F})$
(Lemma \href{https://stacks.math.columbia.edu/tag/03Q8}{lemma higher direct images (Tag 03Q8)})
we obtain a canonical commutative diagram
$$
\xymatrix{
\colim_{i' \geq i}
H^p(U_i \times_{S_i} X_{i'}, \mathcal{F}_{i'}) \ar[r] \ar[d] &
\colim_{i' \geq i}
(R^p g_{i', *}\mathcal{F}_{i'})(U_i \times_{S_i} S_{i'}) \ar[d] \\
H^p(U \times_S X, \mathcal{F}) \ar[r] &
R^pg_*\mathcal{F}(U)
}
$$
Observe that the left hand vertical arrow is an isomorphism
by Theorem \href{https://stacks.math.columbia.edu/tag/09YQ}{theorem colimit (Tag 09YQ)}. We're trying to show that
the right hand vertical arrow is an isomorphism. However, we
already know that the source and target of this arrow
are sheaves on $S_{affine, \etale}$. Hence it suffices to
show: (1) an element in the target, locally comes from an
element in the source and (2) an element in the source
which maps to zero in the target locally vanishes.
Part (1) follows immediately from the above and the fact that
the lower horizontal arrow comes from a map of presheaves
which becomes an isomorphism after sheafification.
For part (2), say $\xi \in  \colim_{i' \geq i}
(R^p g_{i', *}\mathcal{F}_{i'})(U_i \times_{S_i} S_{i'})$
is in the kernel. Choose an $i' \geq i$ and
$\xi_{i'} \in (R^p g_{i', *}\mathcal{F}_{i'})(U_i \times_{S_i} S_{i'})$
representing $\xi$.
Choose a standard \'etale covering
$\{U_{i', k} \to U_i \times_{S_i} S_{i'}\}_{k = 1, \ldots, m}$
such that $\xi_{i'}|_{U_{i', k}}$ comes from
$\xi_{i', k} \in H^p(U_{i', k} \times_{S_{i'}} X_{i'}, \mathcal{F}_{i'})$.
Since it is enough to prove that $\xi$ dies locally, we
may replace $U$ by the members of the \'etale
covering $\{U_{i', k} \times_{S_{i'}} S \to U = U_i \times_{S_i} S\}$.
After this replacement we see that $\xi$ is the image of
an element $\xi'$ of the group
$\colim_{i' \geq i} H^p(U_i \times_{S_i} X_{i'}, \mathcal{F}_{i'})$
in the diagram above. Since $\xi'$ maps to zero in $R^pg_*\mathcal{F}(U)$
we can do another replacement and assume that $\xi'$ maps
to zero in $H^p(U \times_S X, \mathcal{F})$.
However, since the left vertical arrow is an isomorphism
we then conclude $\xi' = 0$ hence $\xi = 0$ as desired.
\end{proof}


\begin{lemma}
\label{lemma-Rf-star-zero-normal-with-alg-closed-function-field}
Let $f : X \to Y$ be a morphism of schemes where $X$ is an integral
normal scheme with separably closed function field. Then
$R^qf_*\underline{M} = 0$ for $q > 0$ and any abelian group $M$.
\end{lemma}

\begin{proof}
Recall that $R^qf_*\underline{M}$ is the sheaf associated
to the presheaf $V \mapsto H^q_\etale(V \times_Y X, M)$ on $Y_\etale$, see
Lemma \href{https://stacks.math.columbia.edu/tag/03Q8}{lemma higher direct images (Tag 03Q8)}.
If $V$ is affine, then $V \times_Y X \to X$ is separated and \'etale.
Hence $V \times_Y X = \coprod U_i$ is a disjoint union of open
subschemes $U_i$ of $X$, see
Lemma \href{https://stacks.math.columbia.edu/tag/09Z9}{lemma normal scheme with alg closed function field (Tag 09Z9)}.
By Lemma \href{https://stacks.math.columbia.edu/tag/09AX}{lemma local rings strictly henselian (Tag 09AX)}
we see that $H^q_\etale(U_i, M)$ is equal to
$H^q_{Zar}(U_i, M)$. This vanishes by
Cohomology, Lemma \href{https://stacks.math.columbia.edu/tag/02UW}{lemma irreducible constant cohomology zero (Tag 02UW)}.
\end{proof}


\begin{lemma}
\label{algebra-lemma-characterize-finite-presentation}
Let $\varphi : R \to S$ be a ring map. The following are equivalent
\begin{enumerate}
\item $\varphi$ is of finite presentation,
\item for every directed system $A_\lambda$ of $R$-algebras
the map
$$
\colim_\lambda \Hom_R(S, A_\lambda) \longrightarrow
\Hom_R(S, \colim_\lambda A_\lambda)
$$
is bijective, and
\item for every directed system $A_\lambda$ of $R$-algebras
the map
$$
\colim_\lambda \Hom_R(S, A_\lambda) \longrightarrow
\Hom_R(S, \colim_\lambda A_\lambda)
$$
is surjective.
\end{enumerate}
\end{lemma}

\begin{proof}
Assume (1) and write $S = R[x_1, \ldots, x_n] / (f_1, \ldots, f_m)$.
Let $A = \colim A_\lambda$. Observe that an $R$-algebra homomorphism
$S \to A$ or $S \to A_\lambda$ is determined by the images of
$x_1, \ldots, x_n$. Hence it is clear that
$\colim_\lambda \Hom_R(S, A_\lambda) \to \Hom_R(S, A)$
is injective. To see that it is surjective, let $\chi : S \to A$
be an $R$-algebra homomorphism. Then each
$x_i$ maps to some element in the image of some $A_{\lambda_i}$.
We may pick $\mu \geq \lambda_i$, $i = 1, \ldots, n$ and
assume $\chi(x_i)$ is the image of $y_i \in A_\mu$ for
$i = 1, \ldots, n$. Consider $z_j = f_j(y_1, \ldots, y_n) \in A_\mu$.
Since $\chi$ is a homomorphism the image of $z_j$ in
$A = \colim_\lambda A_\lambda$ is zero. Hence there exists a
$\mu_j \geq \mu$ such that $z_j$ maps to zero in $A_{\mu_j}$.
Pick $\nu \geq \mu_j$, $j = 1, \ldots, m$. Then the
images of $z_1, \ldots, z_m$ are zero in $A_\nu$. This
exactly means that the $y_i$ map to elements
$y'_i \in A_\nu$ which satisfy the relations $f_j(y'_1, \ldots, y'_n) = 0$.
Thus we obtain a ring map $S \to A_\nu$. This shows that
(1) implies (2).

\medskip\noindent
It is clear that (2) implies (3). Assume (3).
By Lemma \href{https://stacks.math.columbia.edu/tag/00QN}{lemma ring colimit fp (Tag 00QN)} we may write
$S = \colim_\lambda S_\lambda$ with $S_\lambda$
of finite presentation over $R$. Then the identity map
factors as
$$
S \to S_\lambda \to S
$$
for some $\lambda$. This implies that $S$
is finitely presented over $S_\lambda$ by
Lemma \href{https://stacks.math.columbia.edu/tag/00F4}{lemma compose finite type (Tag 00F4)} part (4)
applied to $S \to S_\lambda \to S$. Applying part (2) of the same
lemma to $R \to S_\lambda \to S$ we conclude that $S$ is of finite
presentation over $R$.
\end{proof}


\begin{lemma}
\label{lemma-base-change-field-extension}
Let $L/K$ be an extension of fields. Let $g : T \to S$ be a quasi-compact
and quasi-separated morphism of schemes over $K$. Denote
$g_L : T_L \to S_L$ the base change of $g$ to $\Spec(L)$.
Let $E \in D^+(T_\etale)$ have cohomology sheaves whose stalks
are torsion of orders invertible in $K$. Let $E_L$ be the
pullback of $E$ to $(T_L)_\etale$. Then
$Rg_{L, *}E_L$ is the pullback of $Rg_*E$ to $S_L$.
\end{lemma}

\begin{proof}
If $L/K$ is separable, then $L$ is a filtered colimit of smooth
$K$-algebras, see
Algebra, Lemma \href{https://stacks.math.columbia.edu/tag/07BV}{lemma colimit syntomic (Tag 07BV)}.
Thus the lemma in this case follows immediately from
Lemma \href{https://stacks.math.columbia.edu/tag/0F09}{lemma smooth base change general (Tag 0F09)}.
In the general case, let $K'$ and $L'$ be the perfect closures
(Algebra, Definition \ref{algebra-definition-perfection})
of $K$ and $L$. Then $\Spec(K') \to \Spec(K)$ and
$\Spec(L') \to \Spec(L)$ are universal homeomorphisms
as $K'/K$ and $L'/L$ are purely inseparable
(see Algebra, Lemma \href{https://stacks.math.columbia.edu/tag/0BRA}{lemma p ring map (Tag 0BRA)}).
Thus we have $(T_{K'})_\etale = T_\etale$,
$(S_{K'})_\etale = S_\etale$,
$(T_{L'})_\etale = (T_L)\etale$, and
$(S_{L'})_\etale = (S_L)_\etale$ by
the topological invariance of \'etale cohomology, see
Proposition \href{https://stacks.math.columbia.edu/tag/03SI}{proposition topological invariance (Tag 03SI)}.
This reduces the lemma to the case of the field
extension $L'/K'$ which is separable (by definition of
perfect fields, see Algebra, Definition \ref{algebra-definition-perfect}).
\end{proof}


\begin{lemma}
\label{lemma-smooth-base-change-separably-closed}
Let $K/k$ be an extension of separably closed fields. Let $X$
be a quasi-compact and quasi-separated scheme over $k$.
Let $E \in D^+(X_\etale)$ have cohomology sheaves whose stalks
are torsion of orders invertible in $k$. Then
\begin{enumerate}
\item the maps $H^q_\etale(X, E) \to H^q_\etale(X_K, E|_{X_K})$
are isomorphisms, and
\item $E \to R(X_K \to X)_*E|_{X_K}$ is an isomorphism.
\end{enumerate}
\end{lemma}

\begin{proof}
Proof of (1). First let $\overline{k}$ and $\overline{K}$
be the algebraic closures
of $k$ and $K$. The morphisms $\Spec(\overline{k}) \to \Spec(k)$ and
$\Spec(\overline{K}) \to \Spec(K)$ are universal homeomorphisms
as $\overline{k}/k$ and $\overline{K}/K$ are purely inseparable
(see Algebra, Lemma \href{https://stacks.math.columbia.edu/tag/0BRA}{lemma p ring map (Tag 0BRA)}).
Thus $H^q_\etale(X, \mathcal{F}) =
H^q_\etale(X_{\overline{k}}, \mathcal{F}_{X_{\overline{k}}})$ by
the topological invariance of \'etale cohomology, see
Proposition \href{https://stacks.math.columbia.edu/tag/03SI}{proposition topological invariance (Tag 03SI)}.
Similarly for $X_K$ and $X_{\overline{K}}$.
Thus we may assume $k$ and $K$ are algebraically closed.
In this case $K$ is a limit of smooth $k$-algebras, see
Algebra, Lemma \href{https://stacks.math.columbia.edu/tag/07BV}{lemma colimit syntomic (Tag 07BV)}.
We conclude our lemma is a special case of
Theorem \href{https://stacks.math.columbia.edu/tag/0EYU}{theorem smooth base change (Tag 0EYU)} as reformulated in
Lemma \href{https://stacks.math.columbia.edu/tag/0F09}{lemma smooth base change general (Tag 0F09)}.

\medskip\noindent
Proof of (2). For any quasi-compact and quasi-separated $U$ in $X_\etale$
the above shows that the restriction of the map
$E \to R(X_K \to X)_*E|_{X_K}$ determines an isomorphism on cohomology.
Since every object of $X_\etale$ has an \'etale covering by such $U$
this proves the desired statement.
\end{proof}


\begin{lemma}
\label{lemma-base-change-f-star-general}
Consider the cartesian diagram of schemes
$$
\xymatrix{
X \ar[d]_f & Y \ar[l]^h \ar[d]^e \\
S & T \ar[l]_g
}
$$
Assume that $f$ is flat and every object $U$ of $X_\etale$ has
a covering $\{U_i \to U\}$ such that $U_i \to S$
factors as $U_i \to V_i \to S$ with $V_i \to S$
\'etale and $U_i \to V_i$ quasi-compact with
geometrically connected fibres.
Then for any sheaf $\mathcal{F}$ of sets on $T_\etale$ we have
$f^{-1}g_*\mathcal{F} = h_*e^{-1}\mathcal{F}$.
\end{lemma}

\begin{proof}
Let $U \to X$ be an \'etale morphism such that $U \to S$ factors as
$U \to V \to S$ with $V \to S$ \'etale and $U \to V$ quasi-compact
with geometrically connected fibres. Observe that $U \to V$ is flat
(More on Flatness, Lemma \href{https://stacks.math.columbia.edu/tag/05B9}{lemma etale flat up down (Tag 05B9)}).
We claim that
\begin{align*}
f^{-1}g_*\mathcal{F}(U)
& = g_*\mathcal{F}(V) \\
& = \mathcal{F}(V \times_S T) \\
& = e^{-1}\mathcal{F}(U \times_X Y) \\
& = h_*e^{-1}\mathcal{F}(U)
\end{align*}
Namely, thinking of $U$ as an object of $X_\etale$ and
$V$ as an object of $S_\etale$ we see that the first equality
follows from Lemma \href{https://stacks.math.columbia.edu/tag/0A3H}{lemma sections upstairs (Tag 0A3H)}\footnote{Strictly
speaking, we are also using that the restriction of $f^{-1}g_*\mathcal{F}$
to $U_\etale$ is the pullback via $U \to V$ of the restriction of
$g_*\mathcal{F}$ to $V_\etale$. See
Sites, Lemma \href{https://stacks.math.columbia.edu/tag/03EF}{lemma localize morphism strong (Tag 03EF)}.}.
Thinking of $V \times_S T$ as an object of $T_\etale$
the second equality follows from the definition of $g_*$.
Observe that $U \times_X Y = U \times_S T$ (because $Y = X \times_S T$)
and hence $U \times_X Y \to V \times_S T$
has geometrically connected fibres as a base change of $U \to V$.
Thinking of $U \times_X Y$ as an object of $Y_\etale$, we see that
the third equality  follows from Lemma \href{https://stacks.math.columbia.edu/tag/0A3H}{lemma sections upstairs (Tag 0A3H)}
as before. Finally, the fourth equality follows from the definition
of $h_*$.

\medskip\noindent
Since by assumption every object of $X_\etale$ has an \'etale
covering to which the argument of the previous paragraph applies
we see that the lemma is true.
\end{proof}


\begin{lemma}
\label{lemma-base-change-does-not-hold}
With $f : X \to S$ and $n$ as in Remark \ref{remark-base-change-holds}
assume for some $q \geq 1$ we have that
$BC(f, n, q - 1)$ is true, but $BC(f, n, q)$ is not.
Then there exist a commutative diagram
$$
\xymatrix{
X \ar[d]_f & X' \ar[d] \ar[l] & Y \ar[l]^h \ar[d] \\
S & S' \ar[l] & \Spec(K) \ar[l]
}
$$
with both squares cartesian, where
\begin{enumerate}
\item $S'$ is affine, integral, and normal with algebraically
closed function field,
\item $K$ is algebraically closed and $\Spec(K) \to S'$
is dominant (in other words $K$ is an extension of
the function field of $S'$)
\end{enumerate}
and there exists an integer $d | n$
such that $R^qh_*(\mathbf{Z}/d\mathbf{Z})$ is nonzero.
\end{lemma}

\noindent
Conversely, nonvanishing of $R^qh_*(\mathbf{Z}/d\mathbf{Z})$
in the lemma implies $BC(f, n, q)$ isn't true as
Lemma \ref{lemma-Rf-star-zero-normal-with-alg-closed-function-field}
shows that $R^q(\Spec(K) \to S')_*\mathbf{Z}/d\mathbf{Z} = 0$.

\begin{proof}
First choose a diagram and $\mathcal{F}$ as in
Lemma \ref{lemma-base-change-does-not-hold-pre}.
We may and do assume $S'$ is affine (this is obvious, but
see proof of the lemma in case of doubt).
By Lemma \ref{lemma-base-change-q-integral-top}
we may assume $K$ is algebraically closed.
Then $\mathcal{F}$ corresponds to a $\mathbf{Z}/n\mathbf{Z}$-module.
Such a modules is a direct sum of copies of $\mathbf{Z}/d\mathbf{Z}$
for varying $d | n$ hence we may assume $\mathcal{F}$ is
constant with value $\mathbf{Z}/d\mathbf{Z}$.
By Lemma \ref{lemma-base-change-q-integral-bottom}
we may replace $S'$ by the normalization
of $S'$ in $\Spec(K)$ which finishes the proof.
\end{proof}






\begin{lemma}
\label{lemma-fppf-reduced-fibres-base-change-f-star}
Consider a cartesian diagram of schemes
$$
\xymatrix{
X \ar[d]_f & Y \ar[l]^h \ar[d]^e \\
S & T \ar[l]_g
}
$$
where $f$ is flat and locally of finite presentation
with geometrically reduced fibres.
Then $f^{-1}g_*\mathcal{F} = h_*e^{-1}\mathcal{F}$
for any sheaf $\mathcal{F}$ on $T_\etale$.
\end{lemma}

\begin{proof}
Combine Lemma \ref{lemma-base-change-f-star-general} with
More on Morphisms, Lemma
\href{https://stacks.math.columbia.edu/tag/0EY5}{lemma cover by geometrically connected (Tag 0EY5)}.
\end{proof}


\begin{lemma}
\label{lemma-punctual-base-change}
Let $K$ be a field. For any commutative diagram
$$
\xymatrix{
X \ar[d] & X' \ar[l] \ar[d]_{f'} & Y \ar[l]^h \ar[d]^e \\
\Spec(K) & S' \ar[l] & T \ar[l]_g
}
$$
of schemes over $K$ with
$X' = X \times_{\Spec(K)} S'$ and $Y = X' \times_{S'} T$ and
$g$ quasi-compact and quasi-separated, and every abelian sheaf
$\mathcal{F}$ on $T_\etale$ whose stalks are torsion of orders
invertible in $K$ the base change map
$$
(f')^{-1}Rg_*\mathcal{F}
\longrightarrow
Rh_*e^{-1}\mathcal{F}
$$
is an isomorphism.
\end{lemma}

\begin{proof}
The question is local on $X$, hence we may assume $X$ is affine.
By Limits, Lemma \href{https://stacks.math.columbia.edu/tag/09MV}{lemma relative approximation (Tag 09MV)}
we can write $X = \lim X_i$ as a cofiltered limit with affine
transition morphisms of schemes $X_i$ of finite type over $K$.
Denote $X'_i =  X_i \times_{\Spec(K)} S'$ and $Y_i = X'_i \times_{S'} T$.
By Lemma \ref{lemma-base-change-Rf-star-colim}
it suffices to prove the statement for the squares
with corners $X_i, Y_i, S_i, T_i$.
Thus we may assume $X$ is of finite type over $K$.
Similarly, we may write
$\mathcal{F} = \colim \mathcal{F}[n]$ where the colimit
is over the multiplicative system of integers invertible in $K$.
The same lemma used above reduces us to the case where
$\mathcal{F}$ is a sheaf of $\mathbf{Z}/n\mathbf{Z}$-modules
for some $n$ invertible in $K$.

\medskip\noindent
We may replace $K$ by its algebraic closure $\overline{K}$.
Namely, formation of direct image commutes with base change
to $\overline{K}$ according to
Lemma \ref{lemma-base-change-field-extension} (works for both
$g$ and $h$). And it suffices to prove the agreement after
restriction to $X'_{\overline{K}}$. Next, we may replace
$X$ by its reduction as we have the topological invariance of \'etale
cohomology, see
Proposition \href{https://stacks.math.columbia.edu/tag/03SI}{proposition topological invariance (Tag 03SI)}.
After this replacement the morphism $X \to \Spec(K)$
is flat, finite presentation, with geometrically reduced fibres
and the same is true for any base change, in particular for $X' \to S'$.
Hence $(f')^{-1}g_*\mathcal{F} \to Rh_*e^{-1}\mathcal{F}$
is an isomorphism by Lemma \ref{lemma-fppf-reduced-fibres-base-change-f-star}.

\medskip\noindent
At this point we may apply
Lemma \ref{lemma-base-change-does-not-hold-post}
to see that it suffices to prove: given a commutative diagram
$$
\xymatrix{
X \ar[d]_f & X' \ar[d] \ar[l] & Y \ar[l]^h \ar[d] \\
\Spec(K) & S' \ar[l] & \Spec(L) \ar[l]
}
$$
with both squares cartesian, where $S'$ is affine, integral, and normal
with algebraically closed function field $K$, then
$R^qh_*(\mathbf{Z}/d\mathbf{Z})$ is zero for $q > 0$ and $d | n$.
Observe that this vanishing is equivalent to the statement that
$$
(f')^{-1}R^q(\Spec(L) \to S')_*\mathbf{Z}/d\mathbf{Z}
\longrightarrow
R^qh_*\mathbf{Z}/d\mathbf{Z}
$$
is an isomorphism, because the left hand side is zero for example by
Lemma \ref{lemma-Rf-star-zero-normal-with-alg-closed-function-field}.

\medskip\noindent
Write $S' = \Spec(B)$ so that $L$ is the fraction field of $B$.
Write $B = \bigcup_{i \in I} B_i$ as the union of its finite type
$K$-subalgebras $B_i$. Let $J$ be the set of pairs $(i, g)$ where
$i \in I$ and $g \in B_i$ nonzero with ordering
$(i', g') \geq (i, g)$ if and only if $i' \geq i$ and
$g$ maps to an invertible element of $(B_{i'})_{g'}$.
Then $L = \colim_{(i, g) \in J} (B_i)_g$.
For $j = (i, g) \in J$ set $S_j = \Spec(B_i)$
and $U_j = \Spec((B_i)_g)$.
Then
$$
\vcenter{
\xymatrix{
X' \ar[d] & Y \ar[l]^h \ar[d] \\
S' & \Spec(L) \ar[l]
}
}
\quad\text{is the colimit of}\quad
\vcenter{
\xymatrix{
X \times_K S_j \ar[d] & X \times_K U_j \ar[l]^{h_j} \ar[d] \\
S_j & U_j \ar[l]
}
}
$$
Thus we may apply Lemma \ref{lemma-base-change-Rf-star-colim}
to see that it suffices to prove
base change holds in the diagrams on the right which is what we
proved in Lemma \ref{lemma-kunneth-localize-on-X}.
\end{proof}
\end{document}
